\documentclass[10pt,a4paper]{amsart}
\usepackage[margin=19mm]{geometry}
\usepackage{amsmath,amssymb,mathtools}
\usepackage[scr=rsfs,cal=euler]{mathalfa}
\usepackage{xfrac}
\usepackage[unicode,hidelinks,bookmarks=false]{hyperref}
\newcommand{\ie}{i.e.\ }
\newcommand{\cf}{cf.\ }
\newcommand{\resp}{respectively\ }
\newcommand{\Subsection}{\subsection}
\providecommand{\nobreakdash}{\nobreak\hbox{-}}
\setlength{\emergencystretch}{2em}
\multlinegap0pt
\usepackage{amssymb}
\usepackage{dsfont}
\usepackage{tikz}
\usepackage{mathtools}
\newtheorem{lemma}{Lemma}
\renewcommand{\P}{\mathbb{P}}
\renewcommand{\geq}{\geqslant}
\renewcommand{\leq}{\leqslant}
\def\llbracket{[\hspace{-.10em} [ }
\def\rrbracket{ ] \hspace{-.10em}]}
\title{The extinction of the contact process in a one-dimensional random environment with long-range interactions}
\author{Pablo A. Gomes, Marcelo R. Hil\'ario, Bernardo N. B. de Lima, Thomas Mountford}
\date{Source: arXiv:2506.17444v2 (CC BY 4.0)}
\begin{document}
\maketitle

\noindent\emph{Source and licence.} The extinction of the contact process in a one-dimensional random environment with long-range interactions. Version arXiv:2506.17444v2 (CC BY 4.0), distributed by arXiv under the Creative Commons Attribution 4.0 International licence, http://creativecommons.org/licenses/by/4.0/, as recorded in the arXiv licence metadata for this exact version and shown on the abstract page https://arxiv.org/abs/2506.17444v2. Full licence notice: \texttt{licenses/arxiv-2506.17444-CC-BY-4.0.txt}.\par\smallskip

\noindent\textit{Editorial scope.} Counted unit: the article's lemma l:compare\_site\_bond (source0.tex lines 1261--1273). The article's complete original proof of it is delivered with it (source0.tex lines 1275--1358, one \texttt{proof} environment of 84 lines). No mathematics is written, repaired or re-derived: the statement and the proof are the article's own lines, in the article's own notation, and no retained line is substituted or re-flowed.  Retained uncounted as the source-local context the proof needs: lines 446--454.  Nothing was omitted from any retained span, so the package carries no omission record.  No retained line contains a citation, so the wrapper carries no bibliography and no bibliography entry.  No retained line contains a figure environment, an image-inclusion command or drawing code, so no image is delivered and the package declares no asset dependency.  Editorial formatting only, in addition: an AMS \texttt{amsart} preamble that restates the article's own declarations and macros used by the retained lines, namely the environments \texttt{lemma} on separate counters, together with the standard packages those lines need.  The retained environments print in this document as Lemma 1 in order of appearance, whereas the article prints the same items with the numbering of its own full document.  The complete native source of the article as distributed in the arXiv e-print for this exact version is bundled unchanged as \texttt{sources/180235/source0.tex}.
\begin{equation}
\label{e:def_perc_stret}
\rho_e = \left\{
\begin{array}{ll}
\rho, & \text{if } e = \{(i, j), (i, j+1)\} \text{ for some } i, j, \\
\rho^{\xi_{i+1}}, & \text{if } e = \{(i, j), (i+1, j)\} \text{ for some } i, j.
\end{array}
\right.
\end{equation}

\begin{lemma}
\label{l:compare_site_bond}
Let $\Lambda = \{\ldots, \phi_{-2}+1, \phi_0+1, \phi_2+1, \phi_4+1, \ldots \}$ and $\rho = 1 - (1-p)^{-1/4}$.
Then
\begin{eqnarray}
P_{\Phi,p}\big(\mathcal{C}_{h}( \llbracket 2a+1,2b+1 \rrbracket \times\llbracket c,d \rrbracket )\big)  
& \geq & \P_{\rho}^{\Lambda}\big(\mathcal{C}_{h}(\llbracket a,b \rrbracket \times\llbracket c,d \rrbracket)\big), \\
P_{\Phi,p}\big(\mathcal{C}_{v}( \llbracket 2a+1,2b+1 \rrbracket \times\llbracket c,d \rrbracket )\big), 
& \geq & \P_{\rho}^{\Lambda}\big(\mathcal{C}_{v}(\llbracket a,b \rrbracket \times\llbracket c,d \rrbracket)\big).
\end{eqnarray}


\end{lemma}

\begin{proof}
Let $R = \llbracket 2a+1,2b+1 \rrbracket \times\llbracket c,d \rrbracket $.
Write $e_1 = \langle (0,0), (2,0)\rangle $ and $e_2 = \langle (0,0), (0,1) \rangle$ and define the following set of unoriented edges
\begin{equation*}
e(R) = \big\{ \big\langle (2i+1,j), (2i+1,j) + e_\ell \big\rangle \colon i \in \llbracket a, b\rrbracket, j \in \llbracket c,d \rrbracket, \ell = 1,2 \big\}.
\end{equation*}

Now let $\rho$ be such that $1-p = (1-\rho)^4$ and split each site of the form $(2i+1,j) \in R$  into four sites $(2i+1,j)_1, \dots, (2i+1,j)_4$.
Declare each site $(2i+1,j)_{\ell}$, $\ell = 1, \dots, 4$, open with probability $\rho$, independently.
We can recover the percolation process $P_{\Phi,p}$ inside $R$ by saying that $(2i+1,j)$ is open if, and only if, at least one of its four sites are open.

We now use an exploration process to compare $P_{\Phi,p}(\mathcal{C}_{h}(R))$ with the probability of the same event under the law of the bond percolation on a randomly stretched lattice \eqref{e:def_perc_stret}.

Fix any ordering of the edges in $e(R)$.
Start with 
\begin{equation*}
(A_0,B_0) = \big( \{2a+1\} \times \llbracket c,d \rrbracket, \varnothing \big).
\end{equation*}
Given $(A_n,B_n), n \geq 0$, let
\begin{equation*}
E_n = \{ e \in B_n^c \cap e(R) \colon e = \langle x,y \rangle \textrm{ for some } x \in A_n \textrm{ and } y \notin A_n\}.
\end{equation*}
Inductively, we define  $(A_{n+1},B_{n+1}) = (A_n,B_n)$ if $E_n$ is empty; otherwise, let $e_n = \langle x,y \rangle$ be the earliest edge in $E_n$ with $x = (2i+1, j) \in A_n$.
Set $B_{n+1} = B_n \cup \{e_n\}$, and define $A_{n+1}$ as follows:
\begin{itemize}
\item If $y = x+(2,0)$, set 
\begin{equation}
\label{e:add_y_1}
A_{n+1} =
\begin{cases}
 A_n \cup \{y\}, \textrm{ if $(2i+2, j)$ and $(2i+3,j)_1$ are open,} \\
A_n, \textrm{ otherwise.}
\end{cases}
\end{equation}
\item If $y = x-(2,0)$, set 
\begin{equation}
\label{e:add_y_2}
A_{n+1} =
\begin{cases}
 A_n \cup \{y\}, \textrm{ if $(2i, j)$ and $(2i-1,j)_2$ are open,} \\
A_n, \textrm{ otherwise.}
\end{cases}
\end{equation}
\item If $y = x+(0,1)$, set
\begin{equation}
\label{e:add_y_3}
A_{n+1} =
\begin{cases}
 A_n \cup \{y\}, \textrm{ if $(2i+1, j+1)_3$ is open;,} \\
A_n, \textrm{ otherwise.}
\end{cases}
\end{equation}
\item If $y = x-(0,1)$, set
\begin{equation}
\label{e:add_y_4}
A_{n+1} =
\begin{cases}
 A_n \cup \{y\}, \textrm{ if $(2i+1, j-1)_4$ is open,} \\
A_n, \textrm{ otherwise.}
\end{cases}
\end{equation}
\end{itemize}
At each step of the exploration, the probability of adding vertex $y$ equals either $\rho p^{\phi_{2i}}$, $\rho p^{\phi_{2(i-1)}}$ or $\rho$ depending on whether we are respectively in \eqref{e:add_y_1}, \eqref{e:add_y_2} or in  \eqref{e:add_y_3}, \eqref{e:add_y_4}.

Let $(A_{\infty}, B_{\infty}) = \cup_{n\geq 0}(A_n, B_n)$.
Notice that $A_\infty$ equals the exploration cluster of $A_0$ obtained in the bond percolation model on $e(R)$ where vertical edges are open with probability $\rho$ and horizontal edges $\langle (2i+1,j),(2i+1,j)+(2,0) \rangle$ are open with probability $\rho p^{\phi_{2i}} \geq \rho^{\phi_{2i} + 1}$, $a \leq i < b$.

Therefore,
\begin{align*}
P_{\Phi,p}\big(\mathcal{C}_{h}(R)\big) & 
\geq 
\P\left( A_{\infty}\cap \big(  \{2b+1\} \times \llbracket c,d \rrbracket \big) \neq \varnothing \right) \\ 
% & = \P\left( A_{\infty} \textrm{ has an end-vertex in $\{2b-1\} \times \llbracket c,d \rrbracket$} \right)\\ 
& = \P_{\rho}^{\Lambda}\big(\mathcal{C}_{h}(\llbracket a,b \rrbracket \times\llbracket c,d \rrbracket)\big),
\end{align*}
where, we recall, $\Lambda = \{\ldots, \phi_{-2}+1, \phi_0+1, \phi_2+1, \phi_4+1, \ldots \}$ and $\mathbb{P}^{\Lambda}_{\rho}$ is the quenched law governing the percolation process defined in \eqref{e:def_perc_stret}.
In the right-hand side we abuse notation and write $\mathcal{C}_h \big( (\llbracket a,b \rrbracket \times\llbracket c,d \rrbracket) \big)$ for the event that the respective rectangle is crossed horizontally by a path of  open bonds.
The rectangle $\llbracket a,b \rrbracket \times\llbracket c,d \rrbracket$ appears instead of $R$, because we are using edges of length $2$ in our exploration process.

Similarly
\begin{equation*}
P_{\Phi,p}\big(\mathcal{C}_{v}(R)\big) \geq \P_{\rho}^{\Lambda}\big(\mathcal{C}_{v}(\llbracket a,b \rrbracket \times\llbracket c,d \rrbracket)\big).
\end{equation*}
\end{proof}

\end{document}
